% Generated draft: requires mathematical and source review.
\documentclass[11pt]{article}
\usepackage[margin=25mm]{geometry}
\usepackage{fontspec}
\setmainfont{Latin Modern Roman}
\usepackage{amsmath,amssymb,mathtools,unicode-math}
\setmathfont{Latin Modern Math}
\usepackage{enumitem,needspace,xurl,hyperref}
\setlength{\parindent}{0pt}
\setlength{\parskip}{5pt}
\setlength{\emergencystretch}{3em}
\newcommand{\N}{\mathbb{N}}
\newcommand{\Z}{\mathbb{Z}}
\newcommand{\Q}{\mathbb{Q}}
\newcommand{\R}{\mathbb{R}}
\newcommand{\C}{\mathbb{C}}
\newcommand{\F}{\mathbb{F}}
\newcommand{\A}{\mathbb{A}}
\newcommand{\PP}{\mathbb{P}}
\newcommand{\E}{\mathbb{E}}
\newcommand{\eps}{\varepsilon}
\newcommand{\dd}{\,\mathrm{d}}
\begin{document}
\section*{UnsolvedMath ID: 30005204}
\textbf{Source problem code:} catalogue entry
\subsection{Definitions and mathematical statement}
\textbf{Recoverability qualification.} The supplied record does not determine a complete mathematical conjecture. Its statement begins with an unresolved reference (``this can only be done'') and ends mid-sentence. The registry marks the statement as \texttt{unrecoverable}. Moreover, its supplied source-verification note says that the official report introduces the question as a lead-in and immediately describes it as answered by the reported work. This conflicts with the supplied \texttt{open} label and dated literature assessment. These annotations are preserved as source claims, without independently resolving their conflict.

The intended subject, according to the title and supplied curation, is a classical limit of nonlocal minimal surfaces as \(s\uparrow 1\). The excerpt identifies two requirements: curvature bounds uniform as \(s\uparrow 1\), and a description of interactions between layers that persist in the limit, analogous to a Toda system. It does not specify the nonlocal functional, the meaning or full range of \(s\), the ambient space, the admissible surfaces, the hypotheses for curvature estimates, the topology of convergence, or the limiting conclusion. Thus these requirements cannot be promoted here to a quantified conjecture.

The excerpt compares this proposed passage to the Allen--Cahn limit \(\varepsilon\downarrow 0\), where \(\varepsilon>0\). It denotes an energy by \(J_\varepsilon\), but supplies neither its definition nor the equation, potential, normalization, domain, or boundary conditions. Stability and finite Morse index are invoked for critical points, without specifying the permitted variations. A uniform curvature estimate would require a bound independent of the limiting parameter; the excerpt does not specify the spatial region, constants, or permitted dependence on other data.

As background, the excerpt attributes to Chodosh and Mantoulidis, in its unidentified reference [3], a passage to the Allen--Cahn limit using the finite Morse index of min--max critical points. It describes their work as reproving the existence of infinitely many minimal surfaces in closed three-dimensional manifolds with generic metrics and establishing the multiplicity-one conjecture of Marques and Neves. It identifies two central ingredients: curvature estimates uniform in \(\varepsilon\) for level sets of stable critical points of \(J_\varepsilon\) under an unspecified ``area bound''; and sheet-interaction information involving a Toda system, used to deduce degeneracy of a limiting minimal surface when multiple sheets converge to it. These are supplied historical descriptions, not independently verified theorem statements. Their precise hypotheses, the meaning of degeneracy, and their applicability to the proposed nonlocal limit are not recoverable from this record.
\subsection{Short English statement}
The record concerns curvature control and persistent interactions between layers when nonlocal minimal surfaces approach a classical limit as \(s\uparrow 1\), using the Allen--Cahn limit as an analogy. The excerpt is incomplete and omits the hypotheses and intended conclusion. Its open-status annotation also conflicts with its source-verification note, so no precise open conjecture can be recovered without further source clarification.
\subsection{Sources}
\begin{itemize}
\item Local UnsolvedMath registry, \texttt{data/unsolvedmath/problems.json}, permanent ID \(30005204\), title ``Classical Limit of Nonlocal Minimal Surfaces,'' source problem code \texttt{catalogue entry}; supplied complete source record and its source-verification annotations.
\item Supplied report citation: \emph{Calculus of Variations (2023)}, Oberwolfach Reports, DOI 10.4171/owr/2022/37, pp. 2129--2193. \url{https://doi.org/10.4171/owr/2022/37}. This is also the supplied literature-source link labelled ``original Oberwolfach report.'' The record separately assigns proposed year 2022; the relationship between that year and the citation's 2023 designation is not explained.
\item Supplied literature-source label: ``Wang--Wei, Allen--Cahn transition layers.'' \url{https://doi.org/10.1016/j.aim.2019.106856}. No additional bibliographic details are supplied.
\item The excerpt attributes Allen--Cahn results to Chodosh and Mantoulidis through reference [3], whose bibliographic entry is absent. It also names Marques and Neves in connection with the multiplicity-one conjecture. No proposer is identified for this registry entry.
\item Literature assessment checked at \texttt{2026-08-22}: the supplied dated curation labels the problem open and describes uniform curvature estimates and surviving inter-layer interactions as unresolved requirements. This is supplied dated curation, not independently verified. The supplied source-verification note instead states that the report answers the introductory question and does not present it as open; that conflict remains unresolved.
\end{itemize}
\end{document}
